% GENERATED FILE. Edit canonical lessons, then run npm run generate:books.
% canonical-source-hash: fnv1a64:d99867ff9de35313
% canonical-lessons: ML-C175-snehikkuka, ML-C175-vilkkuka, ML-C175-ayaykkuka, ML-C175-eriyuka, ML-C175-kayaruka

\chapter{Love, Sell, Send, Throw, Climb}
\label{ch:ml-love-sell-send-throw-climb}
\hlchaptermodality{175}

\begin{chapteropening}
\textbf{By the end of this chapter:} \emph{I can say to love, to sell, to send, to throw and to climb.}

Complete the last lesson of chapter 175: I can say to love, to sell, to send, to throw and to climb.
\end{chapteropening}

\section[\texorpdfstring{\ml{സ്നേഹിക്കുക} (snēhikkuka)}{സ്നേഹിക്കുക (snēhikkuka)}]{\ml{സ്നേഹിക്കുക} (snēhikkuka) — to love}

\label{lesson:ML-C175-snehikkuka}



\begin{quote}
Before the new one: say the Malayalam for to find, then the Malayalam for to forget.
\end{quote}

\subsection*{You'll want to know: \ml{സ്നേഹിക്കുക}}
\textbf{\ml{സ്നേഹിക്കുക}} — \emph{snēhikkuka} — \textquotedblleft{}to love\textquotedblright{}.

\begin{grammarlens}[title={What you've built}]
One more word for this chapter.
\end{grammarlens}

\subsection*{Your turn}
\begin{itemize}
\raggedright
  \item \emph{Say it:} \emph{snēhikkuka}
  \item \emph{Say it:} \emph{snēhikkuka}, once more
  \item \emph{Say it:} say \emph{snēhikkuka}
  \item \emph{Recall:} say the Malayalam for to find, then the Malayalam for to forget, then say \emph{snēhikkuka} again
\end{itemize}

\begin{tcolorbox}[breakable,colback=teal!4,colframe=teal!35!black,title={Before you move on}]
What does \ml{സ്നേഹിക്കുക} mean? (\textquotedblleft{}To love\textquotedblright{}.) Say it once more.
\end{tcolorbox}
\section[\texorpdfstring{\ml{വിൽക്കുക} (vilkkuka)}{വിൽക്കുക (vilkkuka)}]{\ml{വിൽക്കുക} (vilkkuka) — to sell}

\label{lesson:ML-C175-vilkkuka}



\begin{quote}
Before the new one: say the Malayalam for to forget, then the Malayalam for to love.

Which marks go around the words somebody said? (Quotation marks.) What can open a line of dialogue? (A dash.)
\end{quote}

\subsection*{You'll want to know: \ml{വിൽക്കുക}}
\textbf{\ml{വിൽക്കുക}} — \emph{vilkkuka} — \textquotedblleft{}to sell\textquotedblright{}.

\begin{grammarlens}[title={What you've built}]
One more word for this chapter.
\end{grammarlens}

\subsection*{Your turn}
\begin{itemize}
\raggedright
  \item \emph{Say it:} \emph{vilkkuka}
  \item \emph{Say it:} \emph{vilkkuka}, once more
  \item \emph{Say it:} say \emph{vilkkuka}
  \item \emph{Recall:} say the Malayalam for to forget, then the Malayalam for to love, then say \emph{vilkkuka} again
\end{itemize}

\begin{tcolorbox}[breakable,colback=teal!4,colframe=teal!35!black,title={Before you move on}]
What does \ml{വിൽക്കുക} mean? (\textquotedblleft{}To sell\textquotedblright{}.) Say it once more.
\end{tcolorbox}
\section[\texorpdfstring{\ml{അയയ്ക്കുക} (ayaykkuka)}{അയയ്ക്കുക (ayaykkuka)}]{\ml{അയയ്ക്കുക} (ayaykkuka) — to send}

\label{lesson:ML-C175-ayaykkuka}



\begin{quote}
Before the new one: say the Malayalam for to love, then the Malayalam for to sell.
\end{quote}

\subsection*{You'll want to know: \ml{അയയ്ക്കുക}}
\textbf{\ml{അയയ്ക്കുക}} — \emph{ayaykkuka} — \textquotedblleft{}to send\textquotedblright{}.

\begin{grammarlens}[title={What you've built}]
One more word for this chapter.
\end{grammarlens}

\subsection*{Your turn}
\begin{itemize}
\raggedright
  \item \emph{Say it:} \emph{ayaykkuka}
  \item \emph{Say it:} \emph{ayaykkuka}, once more
  \item \emph{Say it:} say \emph{ayaykkuka}
  \item \emph{Recall:} say the Malayalam for to love, then the Malayalam for to sell, then say \emph{ayaykkuka} again
\end{itemize}

\begin{tcolorbox}[breakable,colback=teal!4,colframe=teal!35!black,title={Before you move on}]
What does \ml{അയയ്ക്കുക} mean? (\textquotedblleft{}To send\textquotedblright{}.) Say it once more.
\end{tcolorbox}
\section[\texorpdfstring{\ml{എറിയുക} (eṟiyuka)}{എറിയുക (eṟiyuka)}]{\ml{എറിയുക} (eṟiyuka) — to throw}

\label{lesson:ML-C175-eriyuka}



\begin{quote}
Before the new one: say the Malayalam for to sell, then the Malayalam for to send.

Is there a hyphen inside any Malayalam word you have learned? (No.) Which mark sits between alternatives? (The slash.)
\end{quote}

\subsection*{You'll want to know: \ml{എറിയുക}}
\textbf{\ml{എറിയുക}} — \emph{eṟiyuka} — \textquotedblleft{}to throw\textquotedblright{}.

\begin{grammarlens}[title={What you've built}]
One more word for this chapter.
\end{grammarlens}

\subsection*{Your turn}
\begin{itemize}
\raggedright
  \item \emph{Say it:} \emph{eṟiyuka}
  \item \emph{Say it:} \emph{eṟiyuka}, once more
  \item \emph{Say it:} say \emph{eṟiyuka}
  \item \emph{Recall:} say the Malayalam for to sell, then the Malayalam for to send, then say \emph{eṟiyuka} again
\end{itemize}

\begin{tcolorbox}[breakable,colback=teal!4,colframe=teal!35!black,title={Before you move on}]
What does \ml{എറിയുക} mean? (\textquotedblleft{}To throw\textquotedblright{}.) Say it once more.
\end{tcolorbox}
\section[\texorpdfstring{\ml{കയറുക} (kayaṟuka)}{കയറുക (kayaṟuka)}]{\ml{കയറുക} (kayaṟuka) — to climb, to get in}

\label{lesson:ML-C175-kayaruka}



\begin{quote}
Before the new one: say the Malayalam for to send, then the Malayalam for to throw.
\end{quote}

\subsection*{You'll want to know: \ml{കയറുക}}
\textbf{\ml{കയറുക}} — \emph{kayaṟuka} — \textquotedblleft{}to climb, to get in\textquotedblright{}.

\begin{grammarlens}[title={What you've built}]
One more word for this chapter.
\end{grammarlens}

\subsection*{Your turn}
\begin{itemize}
\raggedright
  \item \emph{Say it:} \emph{kayaṟuka}
  \item \emph{Say it:} \emph{kayaṟuka}, once more
  \item \emph{Say it:} say \emph{kayaṟuka}
  \item \emph{Recall:} say the Malayalam for to send, then the Malayalam for to throw, then say \emph{kayaṟuka} again
\end{itemize}

\begin{tcolorbox}[breakable,colback=teal!4,colframe=teal!35!black,title={Before you move on}]
What does \ml{കയറുക} mean? (\textquotedblleft{}To climb, to get in\textquotedblright{}.) Say it once more.
\end{tcolorbox}
